\chapter{哈尔测度　卷积}

希腊人的长度、面积和体积诸概念，本质上都建立在它们在位移下的不变性之上："彼此重合的事物（ἐφαρμόζοντα）相等"（欧几里得《几何原本》第一卷"共同观念"4）；正是通过巧妙地运用这条原理，人们得到了给出经典"图形"（多边形、圆锥曲线、多面体、球面等）面积或体积的全部公式，有时用有限分解的程序，有时用"穷竭法"。¹ 用现代语言来说，希腊几何学家所做的工作，可以视为证明了"集合函数"的存在性——这些函数可加且在位移下不变，但只对十分特殊类型的集合有定义。积分学可以看作是对扩大这类集合函数定义域这一需要的回应，而从卡瓦列里到 H. 勒贝格，正是这一关切处于分析学家研究的首位；至于位移下的不变性，它则沦为次要的考虑，因为这一性质已变成二重或三重积分中变量替换的一般公式以及正交变换行列式等于 ±1 这一事实的平凡推论。即便在非欧几何中（尽管那里的位移群有所不同），观点仍然相同：一般地，黎曼从 \(ds^2\) 出发，用经典的欧氏公式定义面积或体积的无穷小元素（或维数 ≥ 3 时的类似物），而它们在保持 \(ds^2\) 不变的变换下的不变性因此几乎是同义反复。

只是在 1890 年前后，群作用下不变测度概念的其他一些不那么直接的推广，才随着积分不变量理论——尤其是 H. 庞加莱和 E. 嘉当的工作——的发展而出现；

H. 庞加莱只预见到作用于空间一部分上的单参数群，而 E. 嘉当主要关注的是位移群，但它们作用的空间并非其定义所在的空间。例如，他由此（除其他结果外）确定了（{[}52 a{]}, v. II \(_{1}\), pp.~265-302）R \(^{2}\) 或 R \(^{3}\) 中直线空间上的（位移群下不变的）测度；\(^{2}\) 他还指出，李群的积分不变量一般说来无非是特殊的微分不变量，因而有可能用李的方法把它们全部确定出来。然而，在 1897 年 A. 胡尔维茨的那项基本工作{[}168{]}之前，似乎没有人想到在群自身上考虑或使用不变测度。胡尔维茨想要构造（R \(^{n}\) 上的）在正交群下不变的多项式，他从一个评注出发：对有限的线性变换群而言，只要对任意多项式 P 经群中所有元素 s 变换所得的象 s.P 取平均，问题便立即解决；这使他想到，对正交群，可以用关于一个不变测度的积分来代替取平均；他借助欧拉角的参数化表示明确给出了后者的表达式，但随即（独立于 E. 嘉当）注意到，李的方法可以给出所有李群上不变测度的存在性。也许正是由于二十世纪初不变量理论的衰落，胡尔维茨的思想几乎没有引起直接反响，直到 1924 年之后，随着 I. 舒尔和 H. 外尔把弗罗贝尼乌斯关于有限群线性表示的经典理论（第~121 页）推广到紧李群，这些思想才被认为有价值。前者仅限于正交群的情形，并说明胡尔维茨的方法如何能把特征标的经典正交关系推广开来；H. 外尔把这一思想与 E. 嘉当关于半单李代数的工作结合起来，得到了紧李群不可约表示特征标的显式表达式以及完全可约性定理{[}331 b{]}，随后又通过对"正则表示"概念的大胆推广，得到了著名的彼得–外尔定理，它与有限群理论中正则表示分解为其不可约分量的分解完美地类似{[}332{]}。

一年前，O. 施赖尔已经建立了拓扑群的一般理论{[}276 a{]}，从那时起人们已经清楚，彼得–外尔论文中的论证对于任何能在其上定义"不变测度"的拓扑群都将原样成立。说实话，当时拓扑学与测度论的一般概念尚处在全面形成之中，人们有望在其上定义不变测度的是哪一类拓扑群、这种"测度"又要对哪些集合来定义，看来都没有明确的界限。唯一明显的一点是：不能指望把证明李群上不变测度存在性的无穷小方法推广到一般情形。然而，另一股源于勒贝格测度方面工作的思想潮流，恰好引向更直接的求解方法。豪斯多夫在 1914 年证明，在 \(\mathbf{R}^3\) 的全部子集上有定义、在位移下不变且不恒为零的可加集合函数并不存在，于是自然要问这一结果对 \(\mathbf{R}\) 和 \(\mathbf{R}^2\) 是否仍然成立：S. 巴拿赫于 1923 年以一种出人意料的方式解决了这个问题，他证明恰恰相反，在这两种情形这样的"测度"确实存在{[}14 b{]}；他的做法十分巧妙，已经依赖于一个超限归纳构造，并用到函数被群中元素平移后的"平均"\(\frac{1}{n} \sum_{k=1}^{n} f(x + \alpha_k)\)。\(^3\) 正是这类思想使 A. 哈尔得以在 1933 年{[}139{]}迈出决定性的一步：证明具有可数开集基的局部紧群上不变测度的存在性。他受到经典积分学中用棱长任意小的全等立方体的拼接来逼近体积的方法的启发，借助对角线程序，把不变测度作为一列"逼近测度"的极限而得到，这一程序本质上至今仍在使用。这一发现影响巨大，尤其因为它使 J. 冯·诺伊曼能立即对紧群解决希尔伯特著名的"第五问题"，即只凭纯拓扑性质（排除一切预先给定的微分结构）刻画李群的问题。但人们随即注意到，要能有效地使用不变测度，仅知道它存在是不够的，还必须知道它在相差一个常数因子的意义下是唯一的；这一点首先由 J. 冯·诺伊曼对紧群证明，他用的是凭借连续函数的"平均"来定义哈尔测度的方法，与巴拿赫的方法类似{[}324 e{]}；随后 J. 冯·诺伊曼{[}324 f{]}与 A. 韦伊{[}330 c{]}用不同的方法同时得到了局部紧群情形的唯一性，A. 韦伊还同时指明了如何把哈尔的程序推广到一般的局部紧群。也是 A. 韦伊（出处同上）得到了齐性空间上相对不变测度的存在性条件，并最终证明，分离拓扑群上一个"测度"（具有合理的性质）的存在本身就意味着该群是局部预紧的。这些工作实质上完成了哈尔测度的一般理论；此后唯一值得引述的补充是拟不变测度的概念，它直到 1950 年前后才真正出现，与局部紧群在希尔伯特空间中的表示理论相关联。

卷积乘积的历史更为复杂。早在十九世纪初人们就注意到：例如，若 \(F(x,t)\) 是关于 x 与 t 的线性常系数偏微分方程的一个积分，则

\[
\int_ {- \infty} ^ {+ \infty} F (x - s, t) f (s) d s
\]

也是同一方程的积分；泊松等人在 1820 年之前就已经利用这一想法，把热方程的积分写成

\[
\int_ {- \infty} ^ {+ \infty} \exp \left(- \frac {(x - s) ^ {2}}{4 t}\right) f (s) d s.\tag{23.1}
\]

稍后，表达式

\[
\frac {1}{2 \pi} \int_ {- \pi} ^ {+ \pi} \frac {\sin \frac {2 n + 1}{2} (x - t)}{\sin \frac {x - t}{2}} f (t) d t\tag{23.2}
\]

给出傅里叶级数的部分和，而狄利克雷对该积分在 n 趋于 +∞ 时极限的研究，给出了环面 T 上"正则化" \(f \mapsto \rho_{n} * f\) 的第一个例子（实际上是由一列非正"核" \(\rho_{n}\) 作出的，这使它的研究大为复杂）；在"奇异积分"的名义下，类似的积分表达式将成为十九世纪末至二十世纪初的分析学家——从 P. 杜布瓦–雷蒙到 H. 勒贝格——所偏爱的课题。在 R 上，魏尔斯特拉斯利用积分 (23.1) 证明他那条关于多项式逼近的定理，并在此背景下给出了用正"核"序列 \(\rho_{n}\)（形如 \(x \mapsto c_{n} \rho(nx)\)）作正则化的一般原理。在 T 上，用正核作正则化的最著名例子稍后由费耶尔给出，从那时起，这一标准程序将成为绝大多数函数级数"求和法"的基础。

然而，这些工作由于"核"与被正则化的函数二者所扮演的角色缺少对称性，几乎没有显示出卷积乘积的代数性质。对这一点的强调尤其应归功于沃尔泰拉。他一般地研究了两个二元函数的"复合" \(F * G\)

\[
(F * G) (x, y) = \int_ {x} ^ {y} F (x, t) G (t, y) d t
\]

他将其视为两个矩阵的乘积经"从有限过渡到无穷"而得到的推广{[}322 a{]}。他很早便区分出（因其遗传理论中的解释而如此命名的）"闭循环"情形，其中 F 与 G 只依赖于 y - x；此时 \(H = F * G\) 亦然，而且若设 \(F(x, y) = f(y - x)\)，\(G(x, y) = g(y - x)\)，则有

\[
H (x, y) = h (y - x),
\]

其中

\[
h (t) = \int_ {0} ^ {t} f (t - s) g (s) d s
\]

于是，当 \(t \geq 0\) 时，\(h\) 与如下两个函数 \(f_1, g_1\) 的卷积相合：它们分别等于 \(f\) 与 \(g\)（对 \(t \geq 0\)），而对 \(t < 0\) 则等于 0。

与此同时，沃尔泰拉发展起来的代数形式体系并没有揭示出与 R 的群结构以及傅里叶变换的联系。我们在此无须撰写后者的历史；但应当指出，从柯西开始，研究傅里叶积分的分析学家主要关心的是为各种"反演"公式的有效性寻找越来越宽的条件，而多少忽略了它的代数性质。傅里叶本人的工作（以及拉普拉斯关于类似积分 \(\int_{0}^{+\infty} e^{-st} f(t) dt\) 的工作）当然不能这么说；但这些变换本质上是在与线性问题的关联中引入的，因此长久以来人们认为不值得考虑两个傅里叶变换的乘积也就不足为怪了（三角级数或整幂级数的乘积是例外，但它们与离散测度卷积的联系在十九世纪显然不可能被察觉）。这一乘积以及 R 上卷积的最早提及，大概见于切比雪夫的一篇论文{[}306{]}，与概率计算方面的问题有关。事实上，在概率论中，R 上两个"概率律"（总质量为 1 的正测度）的卷积 \(\mu * \nu\) 不是别的，正是 \(\mu\) 与 \(\nu\) 的"乘积"概率律（就相应"随机变量"之和而言）。当然，在切比雪夫那里，涉及的还只是具有（关于勒贝格测度的）密度的概率律的卷积，因而是函数的卷积；况且它在他那里只是偶一出现，在 1920–1930 年之前出现它的那少数几部著作中情形也是如此。1920 年，P.-J. 丹尼尔在一篇很少有人引用的短文{[}76 c{]}中定义了 R 上两个任意测度的卷积以及这种测度的傅里叶变换，并明确指出傅里叶变换把卷积变为通常的乘积；这一形式体系从 1925 年起将被概率学家们大量使用，尤其是追随 P. 莱维。但卷积在群论中的根本重要性直到 1927 年才由 H. 外尔充分认识到；他意识到，对紧群而言，函数的卷积扮演着有限群群代数中乘法的角色，从而使人能据此定义"正则表示"；同时，他还在正则化中找到了有限群代数中单位元的对应物。剩下要做的是把所有这些观点加以综合，这在 A. 韦伊的书中得以完成{[}330 c{]}，它是此后各种推广的前奏：一方面是 I. 盖尔范德的赋范代数，另一方面是广义函数的卷积。

哈尔测度与卷积很快成为不可或缺的工具，体现了如此鲜明地刻画着现代分析的代数化趋向；在后续的历史注记中，我们还须多次提到它们的应用。这里我们只限于指出其中涉及局部紧群闭子群（特别是离散子群）的"变化"的那一项应用。这一理论从数的几何中 K. 马勒的一个结果出发，由 C. 沙博蒂于 1950 年开创，最近则刚被 Macbeath 与 Swierczkowski 相当大地发展和深化了{[}212{]}。

\chapter{非局部紧空间上的积分}

尽管拓扑学与测度论之间联系的研究可以追溯到实变量函数现代理论的开端，非局部紧拓扑空间上的积分却只是到了最近才以一般的方式建立起来。在叙述先于现有综合的那些工作之前，我们将回顾有关拓扑与测度之间关系的观念演变中的几个步骤。

对勒贝格来说，问题只在于对一个或多个实变量的函数求积分。1913 年，拉东在 \(\mathbf{R}^n\) 上定义了一般测度及相应的积分；这一理论在 Ch. de la Vallée Poussin 的著作{[}82{]}中有详细的阐述，并且处处依赖于欧氏空间的拓扑性质。稍后，1915 年，弗雷歇在{[}115 c{]}中在一个配有 \(\sigma\) 代数的集合上定义了"抽象"测度及关于这些测度的积分；他指出，这样便有可能不使用任何拓扑手段而建立勒贝格理论的主要结果。他援引{[}115 d{]}引言第一部分中的如下话语来论证他的工作："例如，在一个具有无穷多坐标的空间中，分析的种种应用已经导出收敛序列的多种互不等价的定义，当其中一种定义被另一种替代时，这些空间中序列族与可加集合函数的性质不会有任何改变。"弗雷歇的研究由卡拉泰奥多里加以完成，集合函数扩张为测度这一重要定理正是后者的贡献。Saks 的书的开头{[}268{]}对这一观点给出了浓缩的阐述。

局部紧群上哈尔测度的发现（第~231 页及以后）及其随即得到的众多应用，随后韦伊与盖尔范德在调和分析方面的工作，导致 1940 年前后这一观点的深刻改变：在这类问题中，最方便的是把测度视为连续函数空间上的一个线性型。这一方法被迫局限于紧空间或局部紧空间，但对几乎全部应用而言并不构成障碍；更好的是，J. 泰特与 A. 韦伊把调和分析引入 p 进群和阿代尔群，使解析数论得到了引人注目的更新。

扩大这一观点、考虑非局部紧拓扑空间上测度的必要性来自完全另一个方向：概率论一步一步地引向这类空间的研究，并提供了大量非平凡的例子。这些发展对测度论的影响之所以姗姗来迟，其原因也许要在概率论的相对孤立中去寻找——直到不久之前，它一直处于传统数学学科的边缘。

\section*{序列空间上的测度}
\phantomsection
\addcontentsline{toc}{section}{序列空间上的测度}

古典概率论中发展最充分的分支之一是极限定理（大数定律、向高斯–拉普拉斯律的趋近，\ldots）；这是对由牵涉极大总体的现象所显示的统计规律性概念的一种深化。这些问题的正确数学表述要求在序列空间上引入测度；这些空间构成有限维空间最明显的推广，是 1920 年前后由弗雷歇、莱维、卢津\ldots{} 开展的"一般分析"研究最钟爱的课题。无论如何，概率论新方法的创立者辛钦与柯尔莫哥洛夫同为卢津的学生，而莱维又很早便转向概率问题，这绝非偶然：这些问题正是新方法的试金石。

序列空间上测度的第一次隐含出现，见于 E. 博雷尔 1909 年论述可数概率的著作{[}32 b{]}。博雷尔一个非常独创的想法，是把他刚刚得到的概率结果应用于证明 0 与 1 之间几乎每个实数的十进展开所具有的性质。这一应用基于下面的基本评注：用每个介于 0 与 1 之间的实数在给定基 \(q\)（\(q \geq 2\)）下展开式中的数字串来定义它；如果随机地依次选取数 \(x\) 的各个数字，彼此独立，并以 \(1/q\) 的等概率取值于 \(0,1,\ldots,q-1\)，则 \(x\) 落入 [0,1] 中某个区间的概率等于该区间的长度。

1923 年，施坦豪斯{[}293{]}严格地建立了这些结果，并描述了博雷尔所考虑的那类无界选择序列的精确数学模型：为简单起见取 \(q = 2\)，并用 \(I\) 记二元集 \(\{0,1\}\)；给 \(I\) 装备测度 \(\mu\)，它由 \(\mu(0) = \mu(1) = \frac{1}{2}\) 定义；乘积空间 \(I^{\mathbf{N}}\) 的元素是取值为 0 或 1 的数构成的序列 \(\varepsilon = (\varepsilon(n))_{n \in \mathbf{N}}\)，而映射 \(\phi : \varepsilon \mapsto \sum_{n \geq 0} \varepsilon(n). 2^{-n-1}\) 除一个可数集外是 \(I^{\mathbf{N}}\) 到区间 \([0,1]\) 上的双射；更有甚者，\(\phi^{-1}\) 把 \([0,1]\) 上的勒贝格测度变为测度 \(P\)，即 \(I^{\mathbf{N}}\) 上各个因子上测度 \(\mu\) 的乘积。事实上，施坦豪斯当时并没有乘积测度的构造；

他利用拟双射 \(\phi\) 的存在性，从 [0,1] 上的勒贝格测度出发构造 \(I^{N}\) 上的测度 P，然后给出 P 的一个公理化刻画。这样得到的同构使人们得以把概率的语言翻译为测度的语言，并应用关于勒贝格积分的已知定理。

在同一篇论文中，施坦豪斯考虑随机级数 \(\sum_{n\geq 0}\sigma_n.a_n\)，其中符号 \(\sigma_{n} = \pm 1\) 是随机选取的，彼此独立且都以同样的概率 \(\frac{1}{2}\) 取值；1928 至 1935 年间，他研究了大量其他的随机级数。另一方面，佩利、维纳与济格蒙德考虑随机傅里叶级数\(^{1}\)，其形如 \(\sum_{-\infty}^{\infty}a_n\exp (2\pi i(nt + \varPhi_n))\)；其中"振幅" \(a_{n}\) 是固定的，而"相位" \(\varPhi_{n}\) 是在 [0, 1] 上均匀分布的独立随机变量。虽然解析上的困难随问题不同而差别极大，但以测度论语言所作的翻译在所有情形都相同，并且代表了博雷尔与施坦豪斯所处理情形的一种推广：问题在于在 \(\mathbf{R}^{\mathbf{N}}\) 上构造一个测度，它是一族测度的乘积，这族测度全都等同于某个质量为 1 的特定正测度 \(\mu\)（在 \(\mathbf{R}\) 上）；例如，上述随机傅里叶级数对应于 \(\mu\) 是 [0, 1] 上勒贝格测度的情形。

为构造这样的乘积测度，可以使用两种方法。第一种是直接方法，由丹尼尔于 1918 年首次建立{[}76 b{]}；1934 年延森重新发现了它{[}173{]}，并对 \(\mu\) 是 \([0,1]\) 上勒贝格测度的情形作了详细研究。第二种方法是寻找与施坦豪斯的技巧类似的技巧，以便化归为 \([0,1]\) 上勒贝格测度的情形；在还没有一般测度论的完整论述可供支配时，这种做法有方便的好处，因为它允许不经新的证明就使用勒贝格的定理。\(^{2}\)

\section*{布朗运动理论}
\phantomsection
\addcontentsline{toc}{section}{布朗运动理论}

这一理论在当代科学的发展中占有特殊的地位，因为它见证了物理问题与"纯"数学之间持续而富有成果的交流。布朗运动由植物学家布朗于 1829 年发现，十九世纪中众多物理学家曾对它进行过密集的研究，\(^{3}\) 但第一个令人满意的数学模型直到 1905 年才由爱因斯坦发明。在粒子沿直线运动的简单情形中，爱因斯坦的基本假设表述如下：若 \(x(t)\) 是时刻 \(t\) 粒子的横坐标，且 \(t_0 < t_1 < \ldots < t_{n-1} < t_n\)，则相继的位移 \(x(t_i) - x(t_{i-1})\)（对 \(1 \leq i \leq n\)）是独立的高斯随机变量。此处不是详细叙述促成爱因斯坦理论的 J. 佩兰那些重要实验工作的场合：就我们的目的而言，只须记住佩兰的一个评注即可：据他说，对布朗运动轨道的观察使他不由自主地想到"数学家们的没有导数的函数"。这一评注正是点燃维纳的最初火花。

另一股完全不同的思想潮流源于气体分子运动论，它由玻尔兹曼与吉布斯在 1870 至 1900 年间发展起来。考虑由 N 个质量为 m 的分子组成、处于（绝对）温度 T 的气体，并用 \(v_{1}, \ldots, v_{N}\) 记该气体 N 个分子的速度；系统的动能等于

\[
\frac {m}{2} (\mathbf {v} _ {1} ^ {2} + \dots + \mathbf {v} _ {N} ^ {2}) = 3 N k T\tag{24.1}
\]

其中 k 是玻尔兹曼常数。按照吉布斯的想法，分子间大量的碰撞使人无法精确确定分子的速度，于是方便的做法是在由方程 (24.1) 定义的 3N 维空间的球面 S 上引入一个概率律 P。"微正则"假设在于假定 P 是球面 S 上旋转不变的质量为 1 的测度。另一方面，麦克斯韦速度定律指出，分子速度分量的概率律是方差为 2kT/m 的高斯测度。E. 博雷尔似乎是最早（1914 年）指出下述事实的人：当分子数很大时，麦克斯韦定律是吉布斯假设与球面性质的一个推论。他考虑高维欧氏空间中的一个球面 S 以及 S 上旋转不变的质量为 1 的测度 P；利用基于斯特林公式的经典逼近方法，他证明 P 在一条坐标轴上的投影近似地为高斯测度。这些结果稍后由加托与莱维加以精确化{[}201 a{]}。给定整数 \(m \geq 1\) 与数 r > 0，用 \(S_{m,r}\) 记形如 \((x_{1}, \ldots, x_{m}, 0, 0, \ldots)\) 且满足 \(x_{1}^{2} + \ldots + x_{m}^{2} = r^{2}\) 的序列所成的集；再用 \(\sigma_{m,r}\) 记 \(S_{m,r}\) 上旋转不变的质量为 1 的测度。用现代语言陈述，加托与莱维的结果如下：测度序列 \(\sigma_{m,1}\) 狭收敛于原点 \((0, 0, \ldots)\) 处的单位质量，而测度序列 \(\sigma_{m,\sqrt{m}}\) 狭收敛于形如下的测度 \(\Gamma\)

\[
d \Gamma (x _ {1}, x _ {2}, \dots) = \prod_ {n = 1} ^ {\infty} d \gamma (x _ {n})
\]

（γ 是 R 上方差为 1 的高斯测度。）

前述测度 \(\Gamma\) 起着无穷维高斯测度的作用。莱维似乎曾朦胧地希望以内在的方式在所有无穷维希尔伯特空间上定义一个高斯测度。事实上，正如莱维与维纳所证明的，测度 \(\Gamma\) 在某种意义下\(^{4}\) 在 \(l^{2}\) 的自同构作用下不变；可惜，集 \(l^{2}\)——即平方可和序列 \((x_{1}, x_{2}, \ldots, x_{n}, \ldots)\) 全体——对 \(\Gamma\) 是零测的。今天我们知道，在无穷维希尔伯特空间上我们只能满足于一个高斯前测度。\(^{5}\)

本质的进展应归功于维纳：既然无穷维希尔伯特空间上不存在合理的高斯测度，那么从高斯前测度出发，借助初等运算便可能在连续函数空间上构造一个测度 w。我们将扼要说明维纳对 w 的最初构造{[}336{]}；它直接受到加托与莱维的关系式 \(\Gamma = \lim_{m \to \infty} \sigma_{m, \sqrt{m}}\) 的影响。对每个整数 \(m \geq 1\)，用 \(H_m\) 记 T = ]0,1] 上在每个区间 \(\left[\frac{k-1}{m}, \frac{k}{m}\right]\)（对 \(k = 1, 2, \ldots, m\)）上取常值的函数所成的集，用 \(\pi_m\) 记 \(R^m\) 中半径为 1 的欧氏球面上旋转不变的质量为 1 的测度。设 \(f_m\) 是 \(H_m\) 到 \(R^m\) 上的同构，它把每个函数所取的值 \(a_k\)（该函数在区间 \(\left[\frac{k-1}{m}, \frac{k}{m}\right]\) 上取常值）对应到向量 \((a_1, a_2 - a_1, \ldots, a_m - a_{m-1})\)（维纳钟爱的"微分空间"一名即由此而来）；用 \(w_m\) 记 \(H_m\) 上的测度，即 \(\pi_m\) 在 \(f_m^{-1}\) 下的象。维纳把所求的测度 w 定义为诸测度 \(w_m\) 的极限。精确地说，用 H 记 T 上带一致收敛拓扑的正则函数所成的集（我们有 \(H_m \subset H\)，对每个整数 \(m \geq 1\)）；对每个在 H 上有界且一致连续的函数 F，极限 \(A\{F\} = \lim_{m \to \infty} \int_{H_m} F(x) dw_m(x)\) 存在；维纳通过对有关各量的涨落作细致的分析得到某些估计，并重新采用丹尼尔提出的紧性论证，证明我们正处在可以应用丹尼尔扩张定理的境地。由此推出存在一个由 \(C(T)\) 承载的测度 w，使得 \(A\{F\} = \int_{C(T)} F(x) dw(x)\)。维纳接着能够证明测度 w 对应于爱因斯坦的假设，\(^{6}\) 而他的估计使他得以给佩兰关于无导数函数的评注以精确的意义：满足 \(\frac{1}{2}\) 阶李普希茨条件的函数所成的集对 w 是可忽略的（相反，对每个满足 \(0 < a < \frac{1}{2}\) 的 a，几乎每个函数都满足 a 阶李普希茨条件）。

今天我们知道维纳测度的许多构造。例如，佩利与维纳使用随机傅里叶级数（{[}243{]}，第~IX 章）：对每个实数序列 \(\mathbf{a} = (a_{n})_{n \geq 1}\) 与每个整数 \(m \geq 0\)，在 ]0, 1] 上定义函数 \(f_{m,a}\) 如下

\[
f _ {m, \mathbf {a}} (t) = a _ {1} t + 2 \sum_ {k = 2} ^ {2 ^ {m + 1}} \frac {1}{\pi k} a _ {k - 1} \sin \pi k t;
\]

可以证明，对 \(\Gamma\)-几乎所有的序列 \(\mathbf{a}\)，函数序列 \(f_{m,\mathbf{a}}\) 趋于一个连续函数 \(f_{\mathbf{a}}\)，且 \(w\) 是 \(\Gamma\) 在（几乎处处有定义的）映射 \(\mathbf{a} \mapsto f_{\mathbf{a}}\) 下的象。后来，莱维在{[}201 b 与 c{]}中给出了另一个构造。最后，卡茨、唐斯克与埃尔德什于 1950 年前后指出，如何在维纳最初的构造中把球面测度 \(\pi_m\)（位于 \(\mathbf{R}^m\) 中）换成更一般的测度。他们的结果在维纳测度与概率论极限定理之间建立了牢固的联系；这些结果将由普罗霍罗夫加以完成和系统化{[}254{]}，对这项工作我们稍后还将回过头来讨论。

这里不是分析由维纳的发现所引起的众多而重要的概率论工作的场合；今天，布朗运动只表现为马尔可夫过程最重要的例子之一。我们只提一下卡茨把维纳测度应用于求解某些抛物型偏微分方程的工作；那是对费曼在量子场论中思想的一种改编——这是数学与物理问题之间相互影响的又一个例子。

\section*{测度的射影极限}
\phantomsection
\addcontentsline{toc}{section}{测度的射影极限}

这是一项特别着眼于概率论的需要而发展起来的理论。关于随机变量的有限序列 \(X_{1}, \ldots, X_{n}\) 的问题，原则上当该序列的律 \(P_{X}\) 已知时便告解决：它是 \(\mathbf{R}^{n}\) 上质量为 1 的一个正测度，使得同时得到诸不等式 \(a_{1} \leq X_{1} \leq b_{1}, \ldots, a_{n} \leq X_{n} \leq b_{n}\) 的概率等于 \(P_{X}(C)\)，其中 \(C\) 是闭长方体 \([a_{1}, b_{1}] \times \ldots \times [a_{n}, b_{n}]\)（位于 \(\mathbf{R}^{n}\) 中）。

实际上，测度 \(P_{X}\) 具有离散支集，或者具有关于勒贝格测度的密度。当涉及随机变量的无穷序列 \((X_{n})_{n\geq1}\) 时，部分序列的律 \(P_{n}\)（即 \((X_{1},\ldots,X_{n})\) 的分布律）一般对每个整数 \(n\geq1\) 都是已知的；这些给定的律满足一个相容性条件，它表达了如下事实：序列 \((P_{n})_{n\geq1}\) 是一个测度的射影子系。直到 1920 年前后，人们还以多少是隐含的方式，通过从有限情形的概率出发作“自然”的过渡到极限，来定义与无穷序列相联系的事件的概率；例如，人们会假定一局博弈终止的概率，是它至多在 n 步内终止的概率当 n 趋于无穷时的极限。自然，这样的理论很难说是自洽的，而且没有什么能排除“悖论”的出现：同一个概率会因用两种各自同样“自然”的程序中的这一种或那一种来计算而得到两个不同的估计。

施坦豪斯{[}293{]}似乎是第一个感到有必要（就掷硬币的情形）不仅考虑射影子系 \((P_{n})_{n\geq1}\) 而且考虑其极限的人。稍早一些，1919 年，丹尼尔{[}76 d{]}已经在一般情形证明了这类射影极限的存在性，\(^{7}\)但这一结果在欧洲似乎一直无人知晓。1933 年，柯尔莫哥洛夫在著作{[}183{]}中重新得到它，他在此书中提出了概率论的公理化概念。丹尼尔与柯尔莫哥洛夫的证明都使用了一个依赖于迪尼定理的紧性论证。

丹尼尔–柯尔莫哥洛夫定理在随机序列 \((X_{n})_{n\geq1}\) 的情形已无可挑剔，但从 1935 年起由柯尔莫哥洛夫、费勒与杜布所开展的对随机函数的研究却隐藏着完全另一量级的困难。例如考虑 R 的一个区间 T，它代表“随机过程”的观测时刻所成的集；可能轨道的全体是乘积空间 \(R^{T}\)，它被看作部分乘积 \(R^{H}\) 的射影极限，其中 H 跑遍 T 的有限子集所成的集；一般就得到一个测度的射影子系 \((\mu_{H})\)。柯尔莫哥洛夫定理确实给出了 \(R^{T}\) 上的一个测度，但它只定义在一个比博雷尔族小得多的族上。\(^{8}\) 柯尔莫哥洛夫构造的一个给出拓扑空间上测度的变体属于角谷{[}177{]}，此后曾被多次重新发现：把 \(\mu_{H}\) 看作 \(\overline{R}^{H}\) 上一个由 \(R^{H}\) 承载的测度，\(^{9}\)紧空间 \(E=\overline{R}^{T}\) 是诸有限乘积 \(\overline{R}^{H}\) 的射影极限，而 E 上的测度 \(\mu\) 可以定义为诸 \(\mu_{H}\) 的射影极限。但这一程序有一个严重的弊病；\(\overline{R}^{T}\) 的元素不具有任何正则性性质，使人能继续该过程的概率研究——甚至不能简单地消去寄生的值 \(\pm\infty\)（由 R 的紧化 \(\overline{R}\) 所引入）。补救的办法是把测度 \(\mu\)（\(\overline{R}^{T}\) 上的）限制到一个特定的子空间（例如布朗运动情形中的 \(C(T)\)）；根本的困难来自这样一个事实：一个函数空间，即使是通常类型的，也未必是 \(\mu\)-可测的（在 \(\overline{R}^{T}\) 中），甚至函数空间的选择本身也可能成问题。\(^{10}\)

1956 年，普罗霍罗夫迈出了决定性的一步，他的著作{[}254{]}对随机过程理论产生了决定性的影响。他把上面分析的文章中维纳所用的方法放进一个方便的公理化形式，从而建立了函数空间上测度的射影极限的一般存在性定理。

1947 年，博赫纳{[}24 b{]}引入了一类较窄的射影子系，它们由有限维实向量空间与满射线性变换组成。这种子系的射影极限以自然的方式等同于代数对偶 \(E^*\)——这里 \(E\) 是装备着弱拓扑 \(\sigma(E^*, E)\) 的一个实向量空间；相应的测度射影子系有一个极限，即测度 \(\mu\)，它只对一个比 \(E^*\) 的博雷尔族小得多的族有定义。博赫纳借其傅里叶变换（它是 \(E\) 上的一个函数）完全刻画了这类“前测度”。但在 \(E\) 上没有拓扑时这一结果几乎无法使用，这时必须考察把 \(\mu\) 看作拓扑对偶 \(E'\)（即 \(E\) 的拓扑对偶）上的测度的可能性。与此独立地，R. 福尔泰与 E. 穆里耶在设法把概率论的某些经典结果（大数定律、中心极限定理）推广到取值于一个巴拿赫空间的随机变量时，也凸现了傅里叶变换在这些问题中所起的基本作用。但实质性的进展直到 1956 年才取得，当时盖尔范德{[}125 c{]}提出，傅里叶变换的自然框架既不是巴拿赫空间也不是希尔伯特空间，而是核型弗雷歇空间。他猜想，这样一个空间上的每个正型连续函数都是其对偶上一个测度的傅里叶变换，这一结果不久便由明洛斯{[}222{]}建立。其重要性首先在于它适用于广义函数空间，而且几乎全部函数空间都是广义函数空间的博雷尔子空间（因而它构成一个比 \(\mathbf{R}^T\) 好得多的框架）。\(^{11}\) 随机广义函数理论是一个正全面展开的领域，我们只须请读者参阅盖尔范德与维连金的著作{[}126{]}即可。

我们刚才提到的关于射影极限的结果都用到了底空间上拓扑的存在性。可以问，在“抽象”测度的情形是否存在一个类似的理论。冯·诺伊曼早在 1935 年就证明了乘积测度在所有情形的存在性，但延森与安德森{[}290{]}发现的一个反例摧毁了每个测度射影子系都承认极限这一希望。人们找到了两种权宜之计：1949 年，C. 扬内斯库-图尔恰在存在适当分解的条件下建立了可数射影极限的存在性，\(^{12}\) 这是研究马尔可夫过程的一个非常有趣的结果；此外，人们已经看清，空间的拓扑只是通过紧子集的全体才介入。因此，很自然地要借助一个集合的子集的紧类这一概念，在一种抽象理论内部把这种状况公理化。这项工作由马尔切夫斯基（他借此建立了一个关于抽象射影极限的定理）与雷尔-纳德泽夫斯基（他研究测度的分解）于 1953 年完成。\(^{13}\)

\section*{一般拓扑空间上的测度与狭收敛}
\phantomsection
\addcontentsline{toc}{section}{一般拓扑空间上的测度与狭收敛}

拓扑学与测度论之间联系的研究尤其被视为对测度的正则性性质的研究，特别是对“外”正则性与“内”正则性的研究；\(^{14}\) 在一个无穷可数的局部紧空间上，内正则性等价于外正则性。勒贝格给出的直线上集合的测度的构造凸显了这两种正则性，而波兰空间上测度的外正则性这一性质，似乎在 1935 年前后已是尽人皆知。但直到 1940 年，A. D. 亚历山德罗夫{[}3{]}才在一篇因战争而传播受阻的文章中凸显了内正则性的作用，并证明波兰空间上的测度具有这种性质；这一结果后来被普罗霍罗夫{[}254{]}重新发现，并且常被错误地归功于这位作者。直到很近的时候人们才看到这一性质可以推广到苏斯林空间；因此，这些空间的重要性大大提高了，尤其是因为人们了解到它们的理论可以不用可度量化假设，而且几乎全部函数空间都是苏斯林的（甚至多数情况下是卢津的）。\(^{15}\)

测度的收敛方式（弱的或狭的）的定义，最方便的做法是把测度空间与一个连续函数空间置于对偶之中。A. A. 马尔可夫推广 F. 里斯的一个古老结果，于 1938 年建立了 \(\mathcal{C}(X)\) 上的正泛函与紧空间 X 上的正则测度之间的一个双射对应。在前面已引用的著作{[}3{]}中，A. D. 亚历山德罗夫把这些结果推广到完全正则空间的情形：他在完全正则空间 X 上的有界连续函数空间 \(\mathcal{C}^{b}(X)\) 上的正线性形式的集合上引入一个层级，\(^{16}\) 他定义了有界测度的狭收敛，并特别证明了下面两个定理：

\begin{enumerate}
\def\labelenumi{\alph{enumi})}
\item
  若 X 是波兰空间，则 \(\mathcal{C}^{b}(X)\) 上与测度相对应的线性形式所成的集在序列的弱收敛下是闭的；
\item
  若有界测度的序列有一个狭极限，则“没有质量逃逸到无穷”（这是普罗霍罗夫狭收敛定理的逆定理的一个弱形式）。
\end{enumerate}

从这些概念与定理的融合中，普罗霍罗夫懂得如何为随机过程理论提取出一些重要结果，并以简单而引人注目的形式呈现它们。在他 1956 年那部已引用过的伟大著作{[}254{]}中，有相当一部分专门讨论波兰空间上的有界正测度；他推广莱维的一个构造，在质量为 1 的正测度所成的集合上定义了一个距离，使它成为一个波兰空间，然后建立了一个关于狭收敛的重要紧性判据。与普罗霍罗夫相独立，勒康{[}197{]}得到了关于测度狭收敛的若干紧性结果；他没有对他所考虑的空间作任何可度量化假设，而且他的结果在局部紧情形化归为迪厄多内先前的定理。

\chapter{李群与李代数}

\section{起源}

将近一个世纪以来被称为“李群论”的这一理论，基本上是由一位数学家建立起来的：索弗斯·李。

在开始叙述它的历史之前，我们将扼要综述为其发展做了准备的若干较早的研究。

\subsection{变换（克莱因–李，1869-1872）。}


1860 年前后，有限集的置换群理论正在发展，并开始得到应用（塞雷、克罗内克、马蒂厄、若尔当）。另一方面，当时正处于全盛时期的不变量理论，使数学家们熟悉了某些对合成封闭的几何变换的无穷集合（特别是线性变换或射影变换）。但是，在若尔当 1868 年关于“运动群”（三维欧几里得空间中位移群的闭子群）的工作{[}174 b{]}之前，这两股思想潮流之间似乎还没有建立起任何自觉的联系。

1869 年，年轻的费利克斯·克莱因（1849-1925），普吕克的学生，在柏林与年长他几岁的挪威人索弗斯·李（1842-1899）结为好友，把他们聚到一起的是两人对普吕克的“直线几何学”、特别是直线丛理论的共同兴趣（第~133 页）。正是在这个时期前后，李形成了他最富原创性的想法之一：把不变量的概念引进分析与微分几何；其来源之一是他观察到，微分方程“用求积法”积分的经典方法全都依赖于这样的事实：方程对于一个“连续”的变换族是不变的。李使用这一想法的第一批工作（由克莱因整理）正是从 1869 年开始的；他在其中研究了“雷耶丛”（在与四面体的各面相交于具有给定交比的 4 个点的直线所成的集）以及以该丛的直线为切线的曲线与曲面（{[}202{]}，vol.~I, Abh. V, pp.~68-77）：他的方法依赖于雷耶丛在保持四面体各顶点不变的三参数交换群（PGL(4, C) 的一个极大环面）下的不变性。同一想法也支配着克莱因与李于 1870 年春同在巴黎时合写的著作（{[}182{]}，v. I, pp.~416-420）；他们在其中实质上确定了平面的射影群 PGL(3, C) 的连通交换子群，并研究了其轨道的几何性质（以 V 曲线或 V 曲面之名）；这使他们通过一个统一的程序，得到各种各样的曲线——代数的或超越的——的性质，例如 \(y = cx^m\) 或对数螺线。两人的证词一致强调伽罗瓦与若尔当的理论给他们留下的深刻印象（若尔当关于伽罗瓦的评注已于 1869 年发表在 Math. Annalen 上；无论如何，李早在 1863 年就已听说过伽罗瓦理论）。1871 年开始对非欧几何感兴趣的克莱因，在其中看到了他为一切已知几何学寻找一个分类原理的研究的开端，这项研究将在 1872 年把他引向“埃尔朗根纲领”。至于李，他在 1873 年给 A. 迈尔的一封信（{[}202{]}，vol.~5, p.~584）中，把他关于变换群的思想的起源定于他旅居巴黎期间，而在 1871 年的一部著作（{[}202{]}，vol.~I, Abh. XII, pp.~153-214）中，他已经使用了“变换群”这一术语，并明确提出了确定 GL(n, C) 的所有子群（“连续的或不连续的”）的问题。说实在的，克莱因与李两人想必都曾难以深入这个新的数学宇宙，克莱因谈到新问世的若尔当的《专题论文集》时，说它是一本“上了七道封印的书”（{[}182{]}，v. I, p.~51）；他在别处谈到（{[}182{]}，v. I, pp.~424-459）时写道：“凡涉及连续算子群的启发式观念的一切，特别是凡涉及微分方程或偏微分方程的积分的一切，都属于李。他后来在他的连续群理论中发展起来的一切概念，当时都已在他那里以萌芽的状态存在，然而却加工得如此之少，以致有许多细节我都不得不说服他，例如甚至在开头，连 V 曲线的存在也是在这样的（说服）之下，即在漫长的交谈过程中（才被接受的）”（{[}182{]}，v. I, p.~415）。

\subsection{无穷小变换}

“无限小”变换的概念可以追溯到无穷小演算的开端；众所周知，笛卡儿通过假设“在无限小之中”每个平面运动都可以等同于一个旋转而发现了瞬时旋转中心；十八世纪解析力学的精心构建完全建立在类似的想法之上。1851 年，西尔维斯特在设法构造线性群

GL(3, C) 或其某些子群的不变量时，给这些矩阵中出现的参数 \(z_{j}\) 以形如 \(\alpha_{j}dt\) 的“无限小”增量，并表达函数 \(f((z_{j}))\) 的不变性，即写下方程 \(f((z_{j} + \alpha_{j}dt)) = f((z_{j}))\)；这给了他一个关于 f 的、带偏导数的线性方程 Xf = 0，其中

\[
X f = \sum_ {j} \alpha_ {j} \frac {\partial f}{\partial z _ {j}},\tag{25.1}
\]

X 因而是一个微分算子，即“一个沿有向参数 \(\alpha_{j}\) 的方向的导数”（{[}304{]}，vol.~3，pp.~326 与 327）；西尔维斯特似乎感觉到其中有一个适用范围相当广的一般原理，但似乎没有再回到这个问题上来。稍后，凯莱（{[}58{]}，v. II, pp.~164-178）对 SL(2, C) 在该群的某些表示中的不变量以同样的方式处理，并证明它满足两个一阶偏微分方程 Xf = 0, Yf = 0，其中 X 与 Y 如上那样从“无限小”变换出发而得到

\[
\left( \begin{array}{c c} 0 & 0 \\ d t & 0 \end{array} \right) \text {   and   } \left( \begin{array}{c c} 0 & d t \\ 0 & 0 \end{array} \right).
\]

用现代的语言，这可以由 X 与 Y 生成李代数 \(\mathfrak{sl}(2,\mathbf{C})\) 这一事实来解释；此外，凯莱还明确地算出了换位积 XY-YX，并证明它也来自一个“无限小”变换。

在若尔当 1868 年关于运动群的论文{[}174 b{]}中，“无限小变换”的概念自始至终被使用，但纯粹是从几何的观点出发。毫无疑问，正是从他那里出现了由一个无限小变换“生成”的单参数群的思想：对若尔当而言，它是通过“适当地重复”该无限小变换而得到的变换全体（loc. cit.，p.~243）。克莱因与李在他们 1871 年的论文中使用了同样的说法“重复的无限小变换”（{[}182{]}，v. I, pp.~424-459），但上下文表明，他们用这个词指的是一个微分方程组的积分。如果他们所考虑的单参数群由变换 \(x' = f(x, y, t)\)，\(y' = g(x, y, t)\) 组成，则相应的“无限小变换”由

\[
d x = p (x, y) d t, d y = q (x, y) d t
\]

其中 \(p(x,y)=\frac{\partial f}{\partial t}(x,y,t_{0}),q(x,y)=\frac{\partial g}{\partial t}(x,y,t_{0})\)，这里 \(t_{0}\) 对应于群的恒等变换。由于克莱因与李显式地知道函数 f 与 g，他们毫不困难地验证了函数

\[
t \mapsto f (x, y, t) \text {   and   } t \mapsto g (x, y, t)
\]

以参数形式给出了微分方程

\[
q (\xi , \eta) d \xi = p (\xi , \eta) d \eta
\]

的经过点 \((x,y)\) 的积分曲线，但对这一点他们没有给出任何一般的理由；此外，在他们的论文的其余部分，他们也没有再用到这一事实。

\subsection{切触变换}

在随后的两年里，李似乎放下了变换群理论（尽管他与克莱因保持着极密切的联系，后者于 1872 年发表了他的“纲领”），转而研究切触变换、一阶偏微分方程的积分以及这两个理论之间的关系。我们不必在此详述这些问题的历史，只须提及几点看来在变换群理论的起源中起了重要作用的内容。

切触变换的概念同时推广了点变换与互反极变换。大体上说，一个切触变换\(^{1}\)（在 \(C^{n}\) 中）是开集 \(\Omega\)（流形 \(T^{\prime}(C^{n})\)——即 \(C^{n}\) 的余切向量所成的流形——的开集）到另一个开集 \(\Omega^{\prime}\)（\(T^{\prime}(C^{n})\) 中的）上的一个同构，它把 \(\Omega\) 的典范 1-形式变为 \(\Omega^{\prime}\) 的典范 1-形式。换句话说，若 \((x_{1},\ldots,x_{n},p_{1},\ldots,p_{n})\) 描述 \(T^{\prime}(C^{n})\) 的典范坐标，则切触变换是一个同构 \((x_{i},p_{i})\mapsto(X_{i},P_{i})\)，它满足关系 \(\sum_{i=1}^{n}P_{i}dX_{i}=\sum_{i=1}^{n}p_{i}dx_{i}\)。这样的变换出现在形如

\[
F \left(x _ {1}, x _ {2}, \dots , x _ {n}, \frac {\partial z}{\partial x _ {1}}, \dots , \frac {\partial z}{\partial x _ {n}}\right) = 0\tag{25.2}
\]

李在研究这些问题的过程中，熟悉了对泊松括号

\[
(f, g) = \sum_ {i = 1} ^ {n} \left(\frac {\partial f}{\partial x _ {i}} \frac {\partial g}{\partial p _ {i}} - \frac {\partial g}{\partial x _ {i}} \frac {\partial f}{\partial p _ {i}}\right)\tag{25.3}
\]

以及类型 (25.1) 的微分算子的括号\(^{2}\) \([X,Y]=XY-YX\)的运用；他把泊松括号 (25.3) 解释为一个与 g 相联系的类型 (25.1) 的变换对 f 的作用，并借此机会观察到，泊松括号的雅可比恒等式意味着与 g 和 h 相对应的微分算子的括号与括号 \((g,h)\) 相联系。求使 \((F,g)=0\) 的函数 g——这出现在雅可比积分偏微分方程 (25.2) 的方法中——对李来说变成了求一个使所给方程保持不变的无穷小切触变换。最后，李被引导去研究函数组 \((u_{j})_{1\leq j\leq m}\)（它们是 \(x_{i}\) 与 \(p_{i}\) 的函数），使得诸括号 \((u_{j},u_{k})\) 是诸 \(u_{h}\) 的函数，他称这些集合为“群”（雅可比实质上已经考虑过它们）。

\section{连续群与无穷小变换}

1873 年秋，李突然重新拾起变换群的研究，并得到了决定性的结果。尽管只能在 1873-1874 年间给 A. 迈尔的几封信（{[}202{]}，vol.~5, pp.~585-608）中追寻他思想展开的脉络，仍可看出他从一个 n 个变量的变换的“连续群”出发

\[
x _ {i} ^ {\prime} = f _ {i} (x _ {1}, \dots , x _ {n}, a _ {1}, \dots , a _ {r}) (1 \leq i \leq n)\tag{25.4}
\]

该群实质上\(^{3}\)依赖 r 个参数 \(a_{1},\ldots,a_{r}\)；他观察到，若变换 (25.4) 当参数取值 \(a_{1}^{0},\ldots,a_{r}^{0}\) 时为恒等变换，\(^{4}\) 则诸 \(x_{i}\) 的限于首阶的泰勒展开式：

\[
f _ {i} \left(x _ {1}, \dots , x _ {n}, a _ {1} ^ {0} + z _ {1}, \dots , a _ {r} ^ {0} + z _ {r}\right) = x _ {i} + \sum_ {k = 1} ^ {r} z _ {k} X _ {k i} \left(x _ {1}, \dots , x _ {n}\right) + \dots (1 \leq i \leq n)\tag{25.5}
\]

给出一个线性地依赖 r 个参数 \(z_{j}\) 的“一般的”无限小变换

\[
d x _ {i} = \left(\sum_ {k = 1} ^ {r} z _ {k} X _ {k i} (x _ {1}, \dots , x _ {n})\right) d t (1 \leq i \leq n).\tag{25.6}
\]

像在他与克莱因合写的论文中那样进行，李积分了微分方程组

\[
\frac {d \xi_ {1}}{\sum_ {k} z _ {k} X _ {k 1} (\xi_ {1} , \dots , \xi_ {n})} = \dots = \frac {d \xi_ {n}}{\sum_ {k} z _ {k} X _ {k n} (\xi_ {1} , \dots , \xi_ {n})} = d t,\tag{25.7}
\]

由此，对每个点 \((z_{1},\ldots,z_{r})\)，他得到一个单参数群

\[
t \mapsto x _ {i} ^ {\prime} = g _ {i} (x _ {1}, \dots , x _ {n}, z _ {1}, \dots , z _ {r}, t) (1 \leq i \leq n)\tag{25.8}
\]

使得对每个 i 有 \(g_{i}(x_{1},\ldots,x_{n},z_{1},\ldots,z_{r},0)=x_{i}\)。他利用变换 (25.4) 组成一个对合成封闭的集合这一事实，以巧妙的方式证明了单参数群 (25.8) 是所给群的一个子群（{[}202{]}，vol.~5, Abh. VII, pp.~32-63）。新的想法——整个理论的关键——是在函数 (25.4) 的泰勒展开中继续推进到二阶。他的论证线索相当混乱且是启发式的（{[}202{]}，vol.~5, Abh. VII, pp.~32-63 以及 {[}202{]}，vol.~5, pp.~600-601）；可以把它整理如下。对于相当小的 \(z_{j}\)，可以在 (25.8) 中取 t=1，这样就得到了群的变换的新参数 \(z_{1},\ldots,z_{r}\)（事实上这是“典范参数”的首次出现）。根据定义，我们关于 (25.7) 有

\[
\frac {\partial g _ {i}}{\partial t} = \sum_ {k} z _ {k} X _ {k i} (x _ {1} ^ {\prime}, \dots , x _ {n} ^ {\prime}),
\]

由此

\[
\frac {\partial^ {2} g _ {i}}{\partial t ^ {2}} = \sum_ {k, j} z _ {k} \frac {\partial X _ {k i}}{\partial x _ {j}} (x _ {1} ^ {\prime}, \dots , x _ {n} ^ {\prime}) \frac {\partial x _ {j} ^ {\prime}}{\partial t} =
\]

\[
\sum_ {k, j} z _ {k} \frac {\partial X _ {k i}}{\partial x _ {j}} (x _ {1} ^ {\prime}, \dots , x _ {n} ^ {\prime}) \left(\sum_ {h} z _ {h} X _ {h j} (x _ {1} ^ {\prime}, \dots , x _ {n} ^ {\prime})\right)
\]

于是得到

\[
\begin{array}{l} x _ {i} ^ {\prime} = x _ {i} + \left(\sum_ {k} z _ {k} X _ {k i} (x _ {1}, \dots , x _ {n})\right) t \\ \quad + \frac {1}{2} \left(\sum_ {k, h, j} z _ {k} z _ {h} \frac {\partial X _ {k i}}{\partial x _ {j}} (x _ {1}, \dots , x _ {n}) X _ {h j} (x _ {1}, \dots , x _ {n})\right) t ^ {2} + \dots , \end{array}
\]

由此，对 \(t = 1\)，得到关于参数 \(z_{j}\) 的泰勒展开式

\[
x _ {i} ^ {\prime} = x _ {i} + \left(\sum_ {k} z _ {k} X _ {k i}\right) + \frac {1}{2} \left(\sum_ {k, h, j} z _ {k} z _ {h} X _ {h j} \frac {\partial X _ {k i}}{\partial x _ {j}}\right) + \dots (1 \leq i \leq n).\tag{25.9}
\]

我们把这些关系简记为向量之间的 \(x' = G(x, z)\)

\[
\boldsymbol {x} = (x _ {1}, \dots , x _ {n}), \boldsymbol {x} ^ {\prime} = (x _ {1} ^ {\prime}, \dots , x _ {n} ^ {\prime}), z = (z _ {1}, \dots , z _ {r});
\]

这些变换的集合对合成封闭这一基本性质写作

\[
G (G (\boldsymbol {x}, \boldsymbol {u}), \boldsymbol {v}) = G (\boldsymbol {x}, H (\boldsymbol {u}, \boldsymbol {v}))\tag{25.10}
\]

其中 \(H = (H_{1},\ldots ,H_{r})\) 与 \(x\) 无关；立即有 \(H(u,0) = u,H(0,v) = v\)，由此得展开式

\[
H _ {i} (u, v) = u _ {i} + v _ {i} + \frac {1}{2} \sum_ {h, k} c _ {i k h} u _ {h} v _ {k} + \dots ,\tag{25.11}
\]

其中未写出的项对 u 或对 v 都是非线性的。借助 (25.9) 与 (25.11) 变换 (25.10)，然后比较两个成员中的 \(u_{h}v_{k}\) 项，李得到关系式

\[
\sum_ {j = 1} ^ {n} \left(X _ {h j} \frac {\partial X _ {k i}}{\partial x _ {j}} - X _ {k j} \frac {\partial X _ {h i}}{\partial x _ {j}}\right) = \sum_ {l = 1} ^ {r} c _ {l h k} X _ {l i} (1 \leq h, k \leq r, 1 \leq i \leq n).\tag{25.12}
\]

他在偏微分方程理论方面的造诣使他把这些条件写成更简单的形式：仿照 (25.1) 的模式，他对这 r 个无限小变换中的每一个——它们由在 (25.6) 中取 \(z_{k}=1\)、\(z_{h}=0\)（对 \(h\neq k\)）而得到——联系上微分算子

\[
A _ {k} (f) = \sum_ {i = 1} ^ {n} X _ {k i} \frac {\partial f}{\partial x _ {i}},\tag{25.13}
\]

并把条件 (25.12) 改写成形式

\[
[ A _ {h}, A _ {k} ] = \sum_ {l} c _ {l h k} A _ {l},\tag{25.14}
\]

这是他的理论的一块基石。在那之前，他不加区分地使用“无限小变换”与“无穷小变换”这两个术语（例如 {[}202{]}，vol.~5, Abh. I, pp.~1-4）；关系式 (25.14) 的简单性使他称算子 (25.13) 为无穷小变换 \(dx_{i} = X_{ki}dt\)（\(1 \leq i \leq n\)）的“符号”（{[}202{]}，vol.~5, Abh. III, pp.~42-75），而且很快，他称为“无穷小变换”的正是算子 (25.13) 本身（{[}202{]}，vol.~5, Abh. III, pp.~42-75 以及 {[}202{]}，vol.~5, p.~589）。

接着他意识到了把“连续群”理论与他先前关于切触变换和偏微分方程的研究联系在一起的紧密纽带。这一联系使他充满热情：“我的旧工作可以说全都预先准备好了，足以为新的变换群理论奠基”，他于 1874 年这样给迈尔写道（{[}202{]}，vol.~5, p.~586）。

在随后的岁月里，李继续研究变换群。除了这里（§ III）综述的一般定理外，他还得到了一些较为特殊的结果：直线与平面上变换群的确定，射影群中小余维数的子群，至多 6 个参数的群，等等。尽管如此，他并没有放弃微分方程。事实上，甚至可以说在他看来，变换群理论本应成为积分微分方程的一种工具，其中变换群所起的作用类似于代数方程的伽罗瓦群所起的作用。\(^{5}\) 我们注意到，这项研究同样引导他引入某些具有无穷多个参数的变换集合，他称之为“无限连续群”；\(^{6}\) 他把“有限连续群”这一名称保留给上面类型 (25.4) 的、具有有限多个参数的变换群。

\section{李群——李代数“辞典”}

李从 1874 年起在大量论文中发展起来的“有限连续”群理论，在宏伟的专著《变换群论》（Theorie der Transformationsgruppen）（{[}203{]}，1888-1893）中得到了系统的阐述，该书是与 F. 恩格尔合作写成的；\(^{7}\) 该理论是第一卷以及第三卷最后五章的对象，第二卷则专用于切触变换。

正如书名所表明的，这部著作中讨论的从来只是变换群，即方程 (25.4) 意义下的变换群，其中“变量”\(x_{i}\) 的空间与“参数”\(a_{j}\) 的空间起初起着同等重要的作用。此外，“抽象”群的概念在当时还没有被清楚地表达出来；1883 年，当李（{[}202{]}，vol.~5, Abh. XII, pp.~311-313）注意到采用 (25.10) 的记号时，方程 \(w = H(u, v)\)——它给出群的两个变换之积的参数——定义了一个新的群，他是把它作为参数空间上的一个变换群来考虑的，由此得到他所说的“参数群”（他甚至得到了两个，它们不是别的，正是左平移群与右平移群）。\(^{8}\)

方程 (25.4) 中的变量 \(x_{i}\) 与参数 \(a_{j}\) 原则上被假定为复的（第三卷第 XIX-XXIV 章除外），函数 \(f_{i}\) 为解析的；李与恩格尔当然意识到，这些函数一般并非对 \(x_{i}\) 与 \(a_{j}\) 的一切复值都有定义，因而这种变换的合成会引起严重的困难（{[}203{]}，v. I, pp.~15-17, pp.~33-40 以及散见各处）；虽然后来他们表述时几乎总是显得他们所研究的变换的合成总是可以无限制地进行，这无疑只是为了叙述的方便，每逢必要他们就明确地重新确立“局部”的观点（例如参见 loc. cit., pp.~168 或 189，或 ibid., v. 3, p.~2, 页末注）；换句话说，他们所研究的数学对象接近于我们今天所说的偏运算律。他们也不忘在适当场合考虑整体群，例如经典群的 4 个系列（{[}203{]}，v. 3, p.~682），但似乎没有给自己提出过如下问题：一般地说一个“整体群”可以是什么；对他们来说，只要能对经典群的“参数”（这些群的“变量”不引起任何困难，因为所涉及的是 \(\mathbf{C}^n\) 的线性变换）得到恒等变换邻域内的“局部”参数系统就够了，而不必关心他们所写下的公式的有效域。然而他们给自己提出了一个明显越出局部理论的问题：⁹ 研究“混合”群，也就是只有有限多个连通分支的群，例如正交群（{[}203{]}，v. I, p.~7）。他们把这项研究表述为对一个对合成和取逆封闭的变换集合的研究，该集合是诸集合 \(H_{j}\) 的并，其中每个集合由形如 (25.4) 的函数组 \(f_{i}^{(j)}\) 所描述；每个 \(H_{j}\) 的（本质）参数的数目甚至被先验地假定依赖于 j，但他们实际上证明这个数目对所有 \(H_{j}\) 都相同。他们的主要结果便是一个有限连续群 G 的存在性，使得 \(H_{j} = G \circ h_{j}\)（对某个 \(h_{j} \in H_{j}\) 以及所有 j）；他们还建立了 G 在混合群中是正规的，并注意到后者的不变量的确定可化归为 G 的不变量与一个不连续群的不变量的确定（{[}203{]}，v. I, chap.~18）。

{[}203{]}中发展起来的一般理论最终（作者们并未以十分系统的方式陈述这一点）锻造出了一部“辞典”，把“有限连续”群的性质翻译为其无穷小变换全体的性质。它建立在“李的三条定理”之上，其中每一条都由一个断言及其逆命题组成。

第一条定理（{[}203{]}，v. I, pp.~33 与 72 以及 v. 3, p.~563）首先断言：若在 (25.4) 中参数是本质的，则函数 \(f_{i}\) 满足形如

\[
\frac {\partial f _ {i}}{\partial a _ {j}} = \sum_ {k = 1} ^ {r} \xi_ {k i} (f (x, a)) \psi_ {k j} (a) (1 \leq i \leq n)\tag{25.15}
\]

其中矩阵 \((\xi_{kl})\) 有最大秩，且 \(\det(\psi_{kj}) \neq 0\)；反过来，若函数 \(f_{i}\) 具有这一性质，则公式 (25.4) 定义一个变换群芽。

第二条定理（{[}203{]}，v. I, pp.~149 与 158，以及 v. 3, p.~590）给出 \(\xi_{ki}\) 之间以及 \(\psi_{ij}\) 之间的关系：关于 \(\xi_{ki}\) 的条件可以写成形式

\[
\sum_ {k = 1} ^ {n} \left(\xi_ {i k} \frac {\partial \xi_ {j l}}{\partial x _ {k}} - \xi_ {j k} \frac {\partial \xi_ {i l}}{\partial x _ {k}}\right) = \sum_ {k = 1} ^ {r} c _ {i j} ^ {k} \xi_ {k l} (1 \leq i, j \leq r, 1 \leq l \leq n)\tag{25.16}
\]

其中 \(c_{ij}^{k}\) 是常数 \((1 \leq i, j, k \leq r)\)，对 i, j 反对称。毛雷尔所给形式的关于 \(\psi_{ij}\) 的条件是：

\[
\frac {\partial \psi_ {k l}}{\partial a _ {m}} - \frac {\partial \psi_ {k m}}{\partial a _ {l}} = \frac {1}{2} \sum_ {1 \leq i, j \leq r} c _ {i j} ^ {k} (\psi_ {i l} \psi_ {j m} - \psi_ {j l} \psi_ {i m}) (1 \leq k, l, m \leq r).\tag{25.17}
\]

引入矩阵 \((\alpha_{ij})\)（它与 \((\psi_{ij})\) 逆步）以及无穷小变换

\[
X _ {k} = \sum_ {i = 1} ^ {n} \xi_ {k i} \frac {\partial}{\partial x _ {i}}, A _ {k} = \sum_ {j = 1} ^ {r} \alpha_ {k j} \frac {\partial}{\partial a _ {j}} (1 \leq k \leq r),\tag{25.18}
\]

(25.16) 与 (25.17) 可以分别写成：

\[
[ X _ {i}, X _ {j} ] = \sum_ {k = 1} ^ {r} c _ {i j} ^ {k} X _ {k} (1 \leq i, j \leq r)\tag{25.19}
\]

\[
\left[ A _ {i}, A _ {j} \right] = \sum_ {k = 1} ^ {r} c _ {i j} ^ {k} A _ {k}.\tag{25.20}
\]

反过来，若给定 r 个无穷小变换 \(X_{k}(1 \leq k \leq r)\)，它们线性无关且满足条件 (25.19)，则由这些变换生成的单参数子群生成一个具有 r 个本质参数的变换群。

最后，第三条定理（{[}203{]}，v. I, pp.~170 与 297 以及 v. 3, p.~597）把满足 (25.19) 的无穷小变换组 \((X_{k})_{1\leq k\leq r}\) 的确定化归为一个纯代数问题：我们必须有

\[
c _ {i j} ^ {k} + c _ {j i} ^ {k} = 0\tag{25.21}
\]

\[
\sum_ {l = 1} ^ {r} (c _ {i l} ^ {m} c _ {j k} ^ {l} + c _ {k l} ^ {m} c _ {i j} ^ {l} + c _ {j l} ^ {m} c _ {k i} ^ {l}) = 0 (1 \leq i, j, k, m \leq r).\tag{25.22}
\]

反过来，\(^{10}\) 若 (25.21) 与 (25.22) 得到满足，就存在一个满足关系式 (25.19) 的无穷小变换组，由此得到一个具有 r 个参数的变换群（换句话说，诸 \(X_{k}\) 的常系数线性组合组成一个李代数，而且反过来每个有限维李代数都能以这种方式得到）。

这些结果由对同构问题的研究加以补充。两个变换群称为相似的，如果可以通过对变量的一个可逆坐标变换和对参数的一个可逆坐标变换从一个变到另一个：从研究之初，李就在“典范参数”的定义中以自然的方式遇到过这一概念。他证明，如果借助“变量”上的一个变换，能把一个群的无穷小变换变成另一个群的无穷小变换，则两群相似（{[}203{]}，v. I, p.~329）。此事成立的一个必要条件是两群的李代数同构，李把这表述为：两群是“gleichzusammengesetzt”（同构组成）的；但这一条件并不充分，整整一章（{[}203{]}，v. I, chap.~19）专门用来求保证两群“相似”的补充条件。另一方面，置换群理论提供了两个这类群的“holoedric 同构”（全同构，即底层“抽象”群的同构）的概念；李把这一概念移植到变换群上，并证明两个这样的群“holoedric 同构”当且仅当它们的李代数同构（{[}203{]}，v. I, p.~418）。特别地，每个变换群都与它的每个参数群全同构，这表明，当人们想研究群的结构时，群所作用的“变量”无关紧要，事实上一切都化归为李代数。\(^{11}\)

同样通过与置换群理论的类比，李引入了子群、正规子群、“meriedric 同构”（满射同态）等概念，并证明它们对应于李代数的子代数、理想与满射同态；此外，他很早就遇到了“meriedric 同构”的一个特别重要的例子，即伴随表示，并且认识了它与群的中心的联系（{[}202{]}，vol.~5, Abh. III, pp.~42-75）。对于这些结果，正如对于基本定理一样，本质的工具是雅可比–克莱布什定理，它给出微分方程组的完全可积性（所谓“弗罗贝尼乌斯定理”的形式之一）；碰巧他还借助单参数群给出了它的一个新证明（{[}203{]}，v. I, chap.~6）。

传递性与本原性的概念——它们对置换群如此重要——对“有限连续”变换群也很自然地出现，李–恩格尔的专著对此作了详细研究（{[}203{]}，v. I, chap.~13 以及散见各处）；与稳定一个点的子群的关系以及齐性空间的概念也被注意到（在不采取整体观点的前提下所能注意到的那部分）（{[}203{]}，v. I, p.~425）。

最后，在 {[}203{]} 中，“辞典”通过引入导群与可解群的概念（李称之为“可积群”；这一由微分方程理论提示的术语一直沿用到 H. 外尔的工作为止）而得以完成（{[}203{]}，v. I, p.~261 以及 v. 3, pp.~678-679）；换位元与括号之间的关系其实李早在 1883 年就已注意到（{[}202{]}，vol.~5, p.~358）。

\subsection*{基本定理的其他证明}
\phantomsection
\addcontentsline{toc}{subsection}{基本定理的其他证明}

在 {[}278 b{]} 中，F. 舒尔证明，在典范坐标中，(15) 的诸 \(\psi_{ik}\) 满足微分方程

\[
\frac {d}{d t} (t \psi_ {i k} (t a)) = \delta_ {i k} + \sum_ {j, l} c _ {j l} ^ {k} t a _ {l} \psi_ {j i} (t a).\tag{25.23}
\]

对这些方程积分，便得到一个等价于公式

\[
\varpi (X) = \sum_ {n \geq 0} \frac {1}{(n + 1) !} (\operatorname{ad} (X)) ^ {n}\tag{25.24}
\]

的公式，它给出指数映射在点 X 处的右微分 \(\varpi(X)\)；特别地，在典范坐标中，诸 \(\psi_{ij}\) 可延拓为诸 \(a_{k}\) 的整函数。F. 舒尔由此导出一个使李先前的一条注记精确化的结果：如果在变换群的定义 (25.4) 中只假设诸 \(f_{i}\) 属于类 \(C^{2}\)，则该群与一个解析群全同构。\(^{12}\) 继他对微分方程组积分的研究之后，E. 嘉当（{[}52 a{]}，v. II₂, p.~371）于 1904 年引入了普法夫型

\[
\omega_ {k} = \sum_ {i = 1} ^ {r} \psi_ {k i} d a _ {i} (1 \leq i \leq r)\tag{25.25}
\]

（用 (25.15) 的记号），后来称为毛雷尔–嘉当形式。于是毛雷尔的条件 (25.17) 可以写成

\[
d \omega_ {k} = - \frac {1}{2} \sum_ {i, j} c _ {i j} ^ {k} \omega_ {i} \wedge \omega_ {j};
\]

E. 嘉当证明，从诸 \(\omega_{k}\) 出发可以展开有限连续群的理论，并确立了这一观点与李的观点的等价性。但对他来说，这个方法的吸引力尤其在于它适用于“无限连续”群——对这些群，他把理论推进得远比利所做的更远——而且它使他得以建立其广义“活动标架”理论。

\section{李代数理论}

一旦变换群与李代数之间的对应建立起来，这一理论就将明显地转向更加代数的方向，并以对李代数的深入研究为中心。\(^{13}\)

从 1888 到 1894 年的第一个短暂时期，以恩格尔、他的学生乌姆劳夫以及特别是基灵与 E. 嘉当的工作为标志，以关于复李代数的一系列引人注目的结果而告终。我们在上面已经看到，可解李代数的概念要归功于李本人，他（在复的情形）证明了可解线性李代数化为三角形约化定理（{[}203{]}，v. I, p.~270）。\(^{14}\) 基灵观察到{[}180{]}，李代数中存在一个最大的可解理想（今天称为根），且李代数对其根的商的根为零；他把根为零的李代数称为半单的，并证明它们是单代数的积（后一概念已经由

李引入，李还证明了“经典”李代数的单性（{[}203{]}，v. 3, p.~682））。

另一方面，基灵在李代数中引入了特征方程 \(\det(\operatorname{ad}(x)-\omega.1)=0\)，李在研究包含李代数中一个给定元素的二维李子代数时已经遇到过它。关于这些方法的分析——基灵以透彻的方式研究半单代数的“一般的”特征方程的根的性质，最终得到他最值得注意的结果，即单（复）李代数的完全确定——我们请读者参看其他历史注记。\(^{15}\)

基灵证明，可解代数的导代数“秩为 0”（这意味着 ad \(x\) 对该代数的每个元素 \(x\) 都是幂零的）。不久之后，恩格尔证明了“秩为 0”的代数是可解的（这一论断实质上就是今天所谓的恩格尔定理）。另一方面，E. 嘉当在他的学位论文中引入了今天所谓的“基灵型”，并建立了借助这个型刻画可解李代数与半单李代数的两个基本判据。

基灵曾断言（{[}180{]}，IV），李代数的导代数是一个半单代数与其幂零的根之和，但他的证明是不完全的。稍后，E. 嘉当未加证明地宣布（{[}52 a{]}，v. I₁, p.~104），更一般地，每个李代数都是其根与一个半单子代数之和；这个方向上当时无可争议地确立的唯一结果，是恩格尔的一个定理，它断言在每个非可解李代数中存在一个三维单李子代数。嘉当论断的第一个公开发表的证明（对复李代数）属于 E. E. 列维{[}200{]}；另一个证明（对实情形同样有效）由 J. H. C. 怀特海于 1936 年给出{[}334 a{]}。1942 年，A. 马尔采夫以关于“列维截口”在共轭意义下的唯一性定理完善了这一结果。

早在最初的工作中，李就给自己提出了每个李代数与一个线性李代数同构的问题。他曾以为通过考虑伴随表示就肯定地解决了它（并由此推出他的“第三定理”的一个证明）（{[}202{]}，vol.~5, Abh. III, pp.~42-75）；他很快认识到他的证明只对中心为零的李代数正确；在他之后，这个问题长时间悬而未决，直到 1935 年才由阿多肯定地解决{[}2 a{]}。另一方面，李实质上还给自己提出过确定单李代数的最小维数线性表示的问题，并对经典代数解决了它；在他的学位论文中，嘉当还对例外单代数解决了这个问题；\(^{16}\) 他在此背景下使用的方法，二十年后将被他推广，以得到实的或复的单李代数的所有不可约表示。

线性表示的完全可约性这一性质，似乎最早是由施图迪（以几何的形式）遇到的。在一份未发表但被（{[}203{]}，v. 3, pp.~785-788）引用的手稿中，他对 SL(2, C) 的李代数的线性表示证明了这一性质，并对 SL(3, C) 与 SL(4, C) 得到了部分结果。李与恩格尔借此机会猜想，完全可约性定理对无论什么 n 的 SL(n, C) 都成立。半单李代数的线性表示的完全可约性由 H. 外尔于 1925 年\(^{17}\)通过一个整体性论证（见后文）建立。第一个代数证明由卡西米尔与范德瓦尔登于 1935 年得到{[}55{]}；其后 R. 布饶尔{[}34 b{]}与 J. H. C. 怀特海{[}334 b{]}给出了另一些代数证明。

最后，在对指数映射的研究（参见后文）中，H. 庞加莱（{[}251{]}，v. III）考虑由一个李代数的算子生成的所有阶微分算子的结合代数；他实质上证明：若 \((X_{i})_{1\leq i\leq n}\) 是李代数的一组基，则由诸 \(X_{i}\) 生成的结合代数以诸 \(X_{i}\) 的某些对称函数（由一个给定单项式通过其因子的所有置换得到的非交换“单项式”之和）为基。他的证明的本质部分是代数性的，使他能够得到包络代数的结构。类似的证明由 G. 伯克霍夫{[}22 b{]}与 E. 维特{[}337 b{]}于 1937 年给出。\(^{18}\)

上面引用的工作大多限于实的或复的李代数，只有它们才对应于通常意义下的李群。对 R 或 C 以外的域上的李代数的研究由雅各布森{[}172 a{]}着手进行，他证明了经典结果的大部分在特征零的域上仍然成立。

\section{指数映射与豪斯多夫公式}

关于指数映射的第一批研究属于 E. 施图迪与 F. 恩格尔；恩格尔{[}104 b{]}注意到，指数映射对 \(\mathbf{SL}(2, \mathbf{C})\) 不是满射（例如 \(\left( \begin{array}{cc} -1 & a \\ 0 & -1 \end{array} \right)\) 当 \(a \neq 0\) 时不是一个指数），但对 \(\mathbf{GL}(n, \mathbf{C})\) 是满射，从而对 \(\mathbf{PGL}(n, \mathbf{C})\) 也是（后一性质已由施图迪对 \(n = 2\) 指出）；于是 \(\mathbf{SL}(2, \mathbf{C})\) 与 \(\mathbf{PGL}(2, \mathbf{C})\) 给出了一个例子：两个群局部同构，但从整体观点看却非常不同。恩格尔还证明，在其余的经典群添上位似之后，指数映射是满射；这项工作被毛雷尔、施图迪等人重新拾起并继续，而没有带来实质性的新结果。

1899 年，H. 庞加莱（{[}251{]}，v. III, pp.~169-172 与 173-212）从一个不同的观点着手指数映射的研究。他的论文似乎写得仓促，因为在几处他断言连通群的每个元素都是指数，而在别处又给出相反的例子。他的结果主要涉及伴随群：他证明，这种群 G 的半单元素可以是李代数 \(L(G)\) 中无穷多个元素的指数，而非半单元素则可能不是指数。若 \(\operatorname{ad}(X)\) 没有作为 \(2\pi i\) 的非零倍数的特征值，则 exp 在 X 处是平展的。他还证明，若 U 与 V 在 \(L(G)\) 中描出回路，且 W 由连续性定义使得 \(e^{U}.e^{V}=e^{W}\)，则我们未必回到 W 的初始判定。他使用一个留数公式，它本质上化归为

\[
\Phi (\operatorname{ad} X) = \frac {1}{2 \pi i} \int \frac {\Phi (\xi) d \xi}{\xi - \operatorname{ad} X}
\]

其中 \(\operatorname{ad}(X)\) 是一个非零特征值均为单重的半单元素，\(\Phi\) 是一个收敛半径足够大的整级数，积分取在包围 ad X 的特征值的一条回路上；他还研究了当 X 趋于一个具有多重特征值的变换时会发生什么。

在庞加莱的工作之前不久，对 \(W\)——作为 \(U\) 与 \(V\) 的函数——在公式 \(e^U.e^V = e^W\) 中的表达式的寻求，就已经是坎贝尔{[}46 a and b{]}两篇论文的主题。稍后贝克写道：“\ldots 李的理论以显然的方式提示，乘积 \(e^Ue^V\) 形如 \(e^W\)，其中 \(W\) 是 \(U\) 与 \(V\ldots\) 的一个交错级数”。这个主题的后续工作旨在把这一断言精确化，并给出 \(W\) 的一个显式公式（或一种构造方法）（“豪斯多夫公式”）。在坎贝尔与庞加莱之后，帕斯卡、贝克{[}13{]}与豪斯多夫{[}152 a{]}又回到这个问题；每个人都认为其前人的证明不能令人信服；主要的困难在于“交错”一词的含义：所涉及的是所考虑的特定李代数的元素，还是普遍的“符号”表达式？坎贝尔、庞加莱与贝克在这一点上都没有表达清楚。相反，豪斯多夫的论文是完全精确的；他首先在有限个不定元的形式结合（非交换）级数代数中工作，并把 U、V、W 看作这个代数的元素。他借助一个与其前人类似的微分方程论证证明了 W 的存在性。同一个论证又被他用来证明当把其中的不定元换成有限维李代数的元素时级数的收敛性。正如贝克所注意到的、庞加莱也独立注意到的，这一结果可以用来给出李第三定理的一个证明；它澄清了群与李代数之间的对应，例如就换位子群而言。

1947 年，邓金{[}98 a{]}重新研究这个问题，并从一开始就考虑赋范李代数（有限维或非有限维，建立在 R、C 或一个超度量域上），从而得到了豪斯多夫公式的显式系数。\(^{19}\)

\section{线性表示与整体李群}

我们所谈论的这些工作，没有一项坦率地触及整体李群的定义与研究问题。沿这条道路的最初几步要追溯到 H. 外尔。他从两个直至当时仍独立发展的理论中得到启发：E. 嘉当的复半单李代数线性表示理论，以及弗罗贝尼乌斯的有限群线性表示理论——后者刚由 I. 舒尔利用胡尔维茨的一个想法移植到正交群上。胡尔维茨曾证明{[}168{]}，通过把有限群上的平均运算换成关于一个不变测度的积分，可以为正交群或酉群构造不变量。他还注意到，把这个方法应用于酉群时，就得到一般线性群的不变量，这是“酉技巧”的第一个例子。1924 年，I. 舒尔{[}279 e{]}用这一程序，通过构造一个正定的非退化不变埃尔米特型，证明了正交群 O(n) 与酉群 U(n) 的表示的完全可约性；他由此借助“酉技巧”推出 O(n, C) 与 SL(n, C) 的全纯表示的完全可约性，为 O(n) 与 U(n) 的特征标建立了正交关系，并确定了 O(n) 的特征标。H. 外尔还把这个方法推广到复半单李代数{[}331 b{]}。给定这样一个代数 g，他证明它具有一个“实紧形式”（这等于说它来自 R 上的一个代数 \(g_{0}\)——其伴随群 \(G_{0}\) 是紧的——把纯量从 R 扩张到 C 而得）。此外他还证明 \(G_{0}\) 的基本群是有限的，因而泛覆盖\(^{20}\)（\(G_{0}\) 的）是紧的。他由此，通过对舒尔程序的一个适当改造，推出 g 的表示的完全可约性，并且还用整体手段给出了 g 的表示的特征标的确定。在一封给 I. 舒尔的信{[}331 a{]}中，H. 外尔概述了舒尔当时尚不知道的嘉当的结果（参见{[}279 e{]}，p.~299, 页脚注），并比较了两种观点：嘉当的方法给出李代数 g 的单连通群的所有全纯表示；在正交群的情形，还得到了一个二叶覆盖（后来称为旋量群）的表示，这是舒尔未曾注意到的；另一方面，舒尔的方法的优点在于能证明完全可约性并显式地给出特征标。

在 H. 外尔的工作之后，E. 嘉当在对对称空间与李群的研究中采取了公开的整体观点。正是这一观点成为他 1930 年对“有限连续”群理论的阐述（{[}52 a{]}，v. I₂, pp.~1165-1225）的基础。其中特别可以找到第三基本定理的整体变体的第一个证明（具有给定李代数的李群的存在性）；嘉当还证明，实李群的每个闭子群都是李群，这推广了 J. 冯·诺伊曼关于线性群的闭子群的一个结果{[}324 b{]}。在这篇论文中，冯·诺伊曼还证明了复半单群的每个连续表示都是实解析的。

在这项工作之后，通常意义（即在 R 或 C 上有限维）下的李群理论其主要轮廓已大体固定。对它的第一个详细阐述由庞特里亚金在他的拓扑群著作{[}253{]}中给出；他在书中坚持一个仍相当接近李的观点，但仔细地区分了局部与整体。随后是谢瓦莱的书{[}62 d{]}，其中还包含对解析流形理论与外微分演算的第一次系统讨论；李的“无穷小变换”在那里表现为向量场，李群 G 的李代数被等同于在 G 下左不变的向量场所成的空间。他把“群芽”方面与“变换群”方面放在了一边。

\section{李群概念的推广}

在今天，李理论的活力既表现于其应用的多样性（在拓扑学、微分几何学、数论等方面），也表现于一些平行理论的创立，在那里底层微分流形的结构被一个相邻的结构（p 进、代数簇、概形、形式概形，\ldots）所代替。我们无须在此提供所有这些发展的历史，只限于其中两种：巴拿赫李群与 p 进 李群。

\subsection{巴拿赫李群}

这里涉及的是“无穷维”李群。从局部观点看，欧几里得空间中 0 的邻域被巴拿赫空间中 0 的邻域所代替。G. 伯克霍夫 1936 年{[}22 a{]}所做的就是这件事，由此得到完备赋范李代数的概念及其与定义在巴拿赫空间的开集上的“群芽”的对应。1950 年前后，邓金把豪斯多夫公式推广到这一情形，完善了这些结果（参见前文）。

伯克霍夫与邓金的定义和结果都是局部的。直到最近，似乎还没有人试图把相应的整体理论明确地建立起来，这无疑是因为缺乏应用。
